\SourceLeaf{212}
Bezeichnungen beibehält, \(\varepsilon\) imaginär, aber \((\frac{1+\varepsilon\cos w}{1-\varepsilon\cos w})^{\frac{1}{2}\varepsilon}\) doch wieder reell wird. Der Vollständigkeit wegen wollen wir jedoch die Formeln für diesen Fall noch besonders beifügen, und gleich Anfangs \(\BareRoot{(\frac{bb}{aa}-1)}=\eta\) setzen. Man hat dann \(w\) durch die Gleichung
\[
\BareRoot{(1+\eta\eta)}\mathord{.}\mathrm{tang}\,u=\mathrm{tang}\,w
\]
\noindent zu bestimmen, und die Differentialgleichung
\[
0=\mathrm{d}t\mp i\,\mathrm{d}w\mathord{.}\frac{1+\eta\eta}{(1+\eta\eta\cos w^2)\sin w}
\]
\noindent wird das Integral
\[
\text{Const.}=t\pm i(\log\mathrm{cotang}\tfrac{1}{2}w+\eta\,\mathrm{Arc}\,\mathrm{tang}.\,\eta\cos w)
\]
\noindent geben, so dass \(X\) für den reellen und \(iY\) für den imaginären Theil von
\[
f(t+i(\log\mathrm{cotang}\tfrac{1}{2}w+\eta\,\mathrm{Arc}\,\mathrm{tang}.\,\eta\cos w))
\]
\noindent wird genommen werden müssen. Die Gegenstücke der beiden obigen speciellen Anwendungen ergeben sich hieraus von selbst. Nach der erstern wird
\[
X=kt,\qquad Y=k\log\mathrm{cotang}\tfrac{1}{2}w+\eta k\,\mathrm{Arc}\,\mathrm{tang}.\,\eta\cos w
\]
\noindent nach der andern
\[
\begin{aligned}
X&=k\,\mathrm{tang}\tfrac{1}{2}w^{\lambda}\mathord{.}e^{-\eta\lambda\,\mathrm{Arc}\,\mathrm{tang}.\,\eta\cos w}\mathord{.}\cos\lambda t\\
Y&=k\,\mathrm{tang}\tfrac{1}{2}w^{\lambda}\mathord{.}e^{-\eta\lambda\,\mathrm{Arc}\,\mathrm{tang}.\,\eta\cos w}\mathord{.}\sin\lambda t
\end{aligned}
\]
\noindent gesetzt werden müssen.

\Article{13}
Als letztes Beispiel wollen wir die allgemeine Darstellung der Oberfläche des Umdrehungs-Ellipsoids auf der Kugelfläche betrachten. Für jenes wollen wir die Bezeichnungen des vorhergehenden Artikels beibehalten, den Halbmesser der Kugelfläche \(=A\), und
\[
\begin{aligned}
X&=A\cos T\sin U\\
Y&=A\sin T\sin U\\
Z&=A\cos U
\end{aligned}
\]
